% ==============================================================================
% 2. METODOLOGIA
% ==============================================================================
\chapter{Metodologia}
\label{cap:metodologia}

\section{Classificação da pesquisa}

Quanto à natureza, esta é uma pesquisa aplicada: o problema é prático (saber o que se recebe ao pedir testes a uma LLM) e o resultado se destina a quem usa essas ferramentas e a quem as avalia. Quanto à abordagem, é quantitativa: todas as variáveis de resposta são medidas por instrumentação, sem julgamento humano na pontuação. Quanto aos objetivos, é explicativa, no sentido de \citeonline{gil2017}, porque isola fatores (modelo, prompt, repetição) e mede o efeito de cada um sobre a distância até o teto. Quanto aos procedimentos, é um experimento controlado, nos termos de \citeonline{wohlin2012}: um objeto fixo, fatores manipulados com níveis definidos, variáveis de resposta medidas da mesma forma em todas as condições e repetições independentes para estimar a variação.

O estudo não parte de uma base de dados existente. Os dados são as métricas de qualidade que a instrumentação produz ao avaliar cada suíte de teste; cada execução de uma suíte sobre o código-alvo rende uma linha de medições.

\section{Objeto, fatores e repetições}

O objeto de estudo é um avaliador de expressões em TypeScript, descrito na seção ``\nameref{sec:alvo}'', fixo em todo o experimento; o Stryker gera para ele 369 mutantes válidos. Há dois fatores manipulados. O fator \textbf{modelo} tem três níveis, escolhidos para cobrir faixas distintas de custo e capacidade: dois modelos de fronteira, um acessado por ferramenta de codificação e outro por aplicativo de chat, e um modelo de baixo custo acessado por API (seção ``\nameref{sec:condicoes}''). O fator \textbf{prompt} tem dois níveis: P1, um pedido básico de testes, e P2, o mesmo pedido acrescido dos critérios pelos quais a suíte será avaliada (Apêndice~\ref{ap:prompts}). Cada combinação modelo~$\times$~prompt é uma \textbf{condição}, e cada condição é gerada $k = 5$ vezes, de forma independente, totalizando 30 suítes. A referência é o \textbf{teto de qualidade}, uma suíte única construída e revisada conforme a seção ``\nameref{sec:teto}''.

A escolha de $k = 5$ é um compromisso entre o custo de gerar e pontuar cada suíte e a necessidade de ter, para cada condição, uma estimativa de dispersão e não só um ponto. A QP4 avalia, a partir da dispersão observada, se esse $k$ é suficiente e quanto seria necessário para uma precisão alvo.

\section{Variáveis de resposta e métricas}
\label{sec:metricas}

Cada suíte é medida nas mesmas seis variáveis, agrupadas em três dimensões:

\begin{enumerate}
    \item \textbf{Cobertura.} Cobertura de linhas (Cob.L) e de ramos (Cob.R), em porcentagem, medidas pelo Vitest com o provedor V8 sobre os três arquivos de lógica do alvo.
    \item \textbf{Poder de detecção de defeitos.} Tratamos cada caso de teste como um detector. O \textit{recall} $R$ é o escore de mutação:
    \begin{equation}
        R = \frac{M_{\text{mortos}} + M_{\text{timeout}}}{M_{\text{válidos}}},
        \label{eq:recall}
    \end{equation}
    em que $M_{\text{válidos}}$ exclui mutantes com erro de compilação ou de execução. A precisão $P$ é a fração dos casos de teste que passam sobre o código correto:
    \begin{equation}
        P = \frac{T - \text{FP}}{T},
        \label{eq:precisao}
    \end{equation}
    em que $T$ é o número de casos de teste e FP é o número deles que falham sobre o código correto (um falso positivo: o teste acusa um defeito que não existe). $F_1$ é a média harmônica de $P$ e $R$.
    \item \textbf{Manutenibilidade.} Densidade de \textit{test smells}: ocorrências detectadas pelos 15 detectores do SNUTS.js \cite{oliveira2024snuts} divididas pelo número de casos de teste.
\end{enumerate}

A variável de interesse é a \textbf{distância até o teto}, calculada por métrica:
\begin{equation}
    \Delta_m = m_{\text{teto}} - m_{\text{suíte}},
    \label{eq:gap}
\end{equation}
positiva quando a suíte fica abaixo do teto (para densidade de \textit{smells}, em que menor é melhor, o sinal é invertido para manter a leitura ``positivo é pior''). Para cada condição reportamos média, desvio padrão, mínimo e máximo das cinco gerações.

Duas decisões sobre as métricas merecem registro. A primeira: o StrykerJS não identifica mutantes equivalentes, então eles permanecem no denominador da \autoref{eq:recall} para \textbf{todas} as suítes, teto incluído; o $R$ reportado é um limite inferior do escore ``verdadeiro'', mas a comparação entre teto e LLMs é feita numa única definição. A segunda: suítes são pontuadas exatamente como geradas, sem correção. Testes que falham sobre o código correto contam como FP e reduzem $P$; apenas na execução de mutação eles são desativados, porque um teste que já falha sobre o código original não é um detector válido e impediria a linha de base verde do Stryker. Corrigir a saída da LLM antes de medi-la contaminaria a medida de qualidade \textit{autônoma}, que é o que se quer medir.

\section{Procedimentos}

O experimento tem quatro etapas. \textbf{(1) Construção e revisão do teto.} A suíte de referência foi construída iterativamente, sem restrição de meios: escrita inicial com apoio de LLM, leitura do relatório de mutação, novos testes para cada mutante sobrevivente, remoção de \textit{smells}, nova rodada, até o escore saturar. Em seguida, cada mutante ainda vivo foi analisado manualmente e classificado como matável (e então morto) ou equivalente, com argumento de por que nenhum teste válido o distingue do original. O critério é o do teste de mutação, não o da API: um mutante em código que nenhuma entrada alcança ainda é matável se um teste puder alcançá-lo por injeção de dependência. O teto só é considerado pronto quando não resta mutante matável, e qualquer suíte gerada que o supere em qualquer métrica reabre a revisão. \textbf{(2) Geração das suítes.} Cada geração parte de uma cópia limpa do projeto-alvo, em diretório neutro, sem histórico, sem os prompts e sem vestígio do aparato de avaliação. A LLM recebe uma única mensagem (o prompt da condição e, para os modelos sem acesso ao sistema de arquivos, o conteúdo dos arquivos do projeto) e devolve o arquivo de testes; não há segunda mensagem, execução de testes com retorno para o modelo nem correção, e a saída é salva como veio. Memória e personalização dos aplicativos ficam desligadas, e cada repetição usa uma sessão nova. \textbf{(3) Pontuação.} Cada suíte é copiada para um ambiente isolado com o código-alvo e pontuada por um único \textit{script}: execução com cobertura e registro de falhas por caso de teste; desativação dos FP; execução do Stryker; contagem de \textit{smells}; cálculo das métricas e das distâncias. Toda a pontuação é \textbf{sequencial}, uma suíte por vez, num mesmo ambiente, pelo motivo descrito na seção ``\nameref{sec:aparato}''. \textbf{(4) Análise.} Para QP1 a QP3, estatística descritiva das distâncias por condição e por dimensão, com intervalos de confiança de 95\% por \textit{bootstrap} percentil (10\,000 reamostragens) sobre as cinco gerações; para o efeito do prompt (QP2), o \textit{delta} de Cliff entre P2 e P1 em cada modelo. Para QP4, desvio padrão e coeficiente de variação entre gerações por condição; a frequência com que o ordenamento das condições pela média se mantém ao tomar uma geração ao acaso por condição; e o número de gerações necessário para que a semi-amplitude do intervalo de 95\% fique em 2 pontos percentuais, por $n = (1{,}96\,s/E)^2$. Com $k = 5$ não há poder para testes de hipótese convencionais, e não os fazemos.

\section{Instrumentos e reprodutibilidade}

O aparato é composto por dois repositórios: o \texttt{bench}, com o código-alvo e uma configuração mínima do Vitest, que é tudo o que uma LLM vê; e o \texttt{validator}, com os prompts, a suíte do teto, as suítes geradas, os \textit{scripts} de geração e de pontuação, a configuração do StrykerJS (8.7, com executor do Vitest 1.6 e cobertura por teste) e os detectores do SNUTS.js. O estudo não envolve participantes humanos. Todos os artefatos, incluindo as respostas completas dos modelos e os resultados em JSON, estão públicos no repositório indicado na apresentação.
